\begin{abstract}
本文检验正且有限历史PB倒数所构成的BP因子，预期高BP对应较高下期收益。采用月末信号、下一交易日收盘换仓、行业与规模控制、加权最小二乘回归（WLS）、等级信息系数（Rank IC）及扣教材成本的等权分层，沪深300、中证500与全部 A 股共同评价2007-02至2026-01归属的228个完整持有期。全部 A 股平均Rank IC为0.0566，正向月份162/228；五层年化净收益由4.55\%递增至20.98\%，多空年化13.00\%，Q5相对内部基准的算术年化超额5.31\%。指数池方向相同而强度较弱。行业共同窗口中Q5在19/30个可比行业排名第一；房地产、钢铁等行业为负超额，综合行业样本不足。2025年正IC与负多头超额并存，低市值域的排序关系更强。结论是保留BP作为后续研究候选，观察时变性与行业边界；历史版本修订、估值近似及简化成本使本结果不能直接解释为严格PIT或实盘可实现收益。
\end{abstract}
\noindent\textbf{关键词：}账面市值比；单因子测试；Rank IC；行业中性化；分层回测。
